% =====================================================================
%  1bat-ud6-3.tex  —  SESSIÓ 3: Consolidació: probabilitats simples i les seves formes
%  (sense preàmbul: s'inclou des de main.tex)
% =====================================================================
\horatitol{Sessió 3 — Consolidació: probabilitats simples i les seves formes}%
{Abans dels successos compostos, assentem el càlcul de probabilitats simples caçant els errors típics i movent-nos amb soltesa entre fracció, decimal i percentatge.}

\subsection{Idea clau}

Ja sabem aplicar la regla de Laplace (sessió 2). Abans d'avançar cap als successos
\textbf{compostos} i \textbf{condicionats} —la part difícil de la unitat—, avui
\textbf{consolidem} el càlcul de probabilitats simples i la conversió entre les seves
tres formes. No hi ha matèria nova: es tracta de fer-ho \textbf{bé} i de caçar els
errors de sempre.

\begin{idea}
\textbf{Els tres errors que cacem avui}
\begin{enumerate}[leftmargin=1.6em, itemsep=2pt]
  \item \textbf{Comptar malament els casos possibles} (oblidar-ne algun, o comptar-ne de
        més).
  \item \textbf{Donar una probabilitat més gran que $1$} (o negativa): senyal segur
        d'error.
  \item \textbf{Equivocar-se en passar de fracció a percentatge} (recorda: divideix i
        multiplica per $100$).
\end{enumerate}
\textbf{Comprova sempre:} tota probabilitat ha de quedar \textbf{entre $0$ i $1$}, i la
suma de les probabilitats de tots els resultats possibles ha de fer \textbf{$1$}.
\end{idea}

\subsection{Exemples}

\begin{exemple}[caça l'error: probabilitat més gran que 1]
Algú calcula la probabilitat de treure un múltiple de $2$ en un dau així: «hi ha $3$
parells, i en surten $2$ de cada cara, $\frac{6}{6}\cdot\dots$» i acaba donant
$P=1,5$. \textbf{És impossible:} cap probabilitat pot ser més gran que $1$. El càlcul
correcte: parells $\{2,4,6\}$, $3$ casos de $6$:
\[
P=\frac{3}{6}=\frac{1}{2}=0,5.
\]
\textbf{Correcció:} $0,5$, no $1,5$.
\end{exemple}

\begin{exemple}[caça l'error: passar a percentatge]
Algú diu que $\frac{1}{4}$ és «el $4\%$». \textbf{És incorrecte:} per passar a
percentatge cal dividir i multiplicar per $100$:
\[
\frac{1}{4}=0,25=25\%.
\]
\textbf{Correcció:} $\frac{1}{4}$ és el $25\%$, no el $4\%$. (Confondre el denominador
amb el percentatge és un error molt comú.)
\end{exemple}

\subsection{Exercicis resolts}

\begin{exresolt}[probabilitat d'un perfil i el seu contrari]
En un grup de $40$ usuaris d'un servei, $24$ tenen estudis bàsics, $12$ estudis mitjans
i $4$ estudis superiors. Si en triem un a l'atzar, calcula la probabilitat de cada
perfil (en les tres formes) i comprova que sumen $1$.
\solucio
\[
P(\text{bàsics})=\frac{24}{40}=\frac{3}{5}=0,6=60\%,
\]
\[
P(\text{mitjans})=\frac{12}{40}=\frac{3}{10}=0,3=30\%,
\]
\[
P(\text{superiors})=\frac{4}{40}=\frac{1}{10}=0,1=10\%.
\]
\textbf{Comprovació:} $0,6+0,3+0,1=1$ (i $60\%+30\%+10\%=100\%$): correcte.
\resposta{$60\%$, $30\%$, $10\%$; sumen $100\%$.}
\end{exresolt}

\begin{exresolt}[del percentatge a la fracció]
Es diu que la probabilitat que un usuari assisteixi a una activitat és del $75\%$.
Expressa-la en decimal i en fracció, i calcula la probabilitat que \textbf{no} hi
assisteixi.
\solucio
$75\%=0,75=\dfrac{75}{100}=\dfrac{3}{4}$.

Probabilitat de no assistir (succés contrari):
\[
1-0,75=0,25=\frac{1}{4}=25\%.
\]
\resposta{$0,75=\frac{3}{4}$; no assistir $=0,25=\frac{1}{4}=25\%$.}
\end{exresolt}

\subsection{Exercicis per fer}

\noindent\textit{Itinerari: els exercicis 1–4 són la base; el 5 i el 6 són ampliació
(successos contraris i grups d'usuaris).}

\begin{exercicis}
\begin{exlist}
  \item \textbf{(Caça l'error)} En cada cas, digues si és correcte i, si no, corregeix-lo:
    \begin{enumerate}[label=\alph*), leftmargin=1.6em, topsep=2pt]
      \item «La probabilitat de treure un $4$ en un dau és $\frac{4}{6}$.»
      \item «$\frac{1}{2}$ és el $2\%$.»
      \item «La probabilitat d'un succés és $1,2$.»
    \end{enumerate}
  \item Calcula i expressa en les tres formes la probabilitat de treure, d'una baralla
        de $40$ cartes: (a) un bastó (n'hi ha $10$); (b) una figura (sota, cavall i rei
        de cada pal: $12$ en total).
  \item En una bossa amb $8$ boles vermelles, $6$ de blaves i $6$ de verdes, calcula la
        probabilitat de cada color i comprova que sumen $1$.
  \item Passa a percentatge: (a) $\frac{2}{5}$; (b) $0,08$; (c) $\frac{7}{20}$.
  \item \textbf{(Ampliació — contrari)} Si la probabilitat que plogui demà és del $30\%$,
        quina és la probabilitat que \textbf{no} plogui? Expressa-la en les tres formes.
  \item \textbf{(Ampliació — usuaris)} En un casal de $60$ infants, $45$ van a primària i
        la resta a secundària. Calcula la probabilitat que un infant triat a l'atzar
        vagi a secundària, en fracció, decimal i percentatge.
\end{exlist}
\end{exercicis}

% ---------------------------------------------------------------------
\begin{idea}
\textbf{Full d'autoavaluació} — marca amb sinceritat com et trobes, de cara a la prova:
\begin{center}
\renewcommand{\arraystretch}{1.4}
\begin{tabular}{p{8.2cm} c c c}
\toprule
\textbf{Sé fer\dots} & \textbf{Sol} & \textbf{Amb ajuda} & \textbf{Encara no} \\
\midrule
Comptar bé casos favorables i possibles & \rule{0.7cm}{0pt} & & \\
Aplicar la regla de Laplace & & & \\
Passar entre fracció, decimal i percentatge & & & \\
Calcular el succés contrari & & & \\
Comprovar que la probabilitat és entre $0$ i $1$ & & & \\
\bottomrule
\end{tabular}
\end{center}
\end{idea}

\noindent\textbf{Atenció a la diversitat.} Si encara et costa, comença per l'\emph{itinerari
mínim} (exercicis amb casos fàcilment comptables i conversió guiada entre formes) i
tingues a mà la fitxa de suport de la sessió 1; si ja domines el càlcul, fes les
\emph{ampliacions} (exercicis 5 i 6, amb contraris i grups d'usuaris).

% ---------------------------------------------------------------------
\fullsolucions{Sessió 3}
\begin{solucions}
\begin{sollist}
  \item (a) \textbf{Incorrecte}: només hi ha \emph{un} $4$, de $6$ cares;
        $P=\frac{1}{6}$. \quad (b) \textbf{Incorrecte}: $\frac{1}{2}=0,5=50\%$, no el
        $2\%$. \quad (c) \textbf{Incorrecte}: cap probabilitat pot superar $1$.
  \item (a) $\frac{10}{40}=\frac{1}{4}=0,25=25\%$; (b)
        $\frac{12}{40}=\frac{3}{10}=0,3=30\%$.
  \item Total $8+6+6=20$. Vermella $\frac{8}{20}=0,4=40\%$; blava
        $\frac{6}{20}=0,3=30\%$; verda $\frac{6}{20}=0,3=30\%$. Sumen $100\%$.
  \item (a) $\frac{2}{5}=0,4=40\%$; (b) $0,08=8\%$; (c) $\frac{7}{20}=0,35=35\%$.
  \item No plogui: $1-0,30=0,70=\frac{7}{10}=70\%$.
  \item Secundària: $60-45=15$ infants. $P=\frac{15}{60}=\frac{1}{4}=0,25=25\%$.
\end{sollist}
\end{solucions}
